


Our goal is to develop a system that assists LLMs in proof generation.
As a proof of concept, we focus on geometry problems from the FormalGeo-7k dataset~\cite{zhang2023formalgeo}, containing 6,981 SAT-level Euclidean geometry problems.
% \footnote{Available at \url{https://github.com/FormalGeo/FormalGeo} under the MIT License. Used here solely for non-commercial academic research.}.


%The dataset includes natural language descriptions, geometric diagrams, formal representations, and annotated proof sequences 

Our approach is two-fold (see Figure~\ref{fig:pipeline}). Given a target problem, we first %identify {\bf analogous problems and their corresponding proofs}
%This approach 
build on the insight that structurally similar problems often admit similar proofs. We abstract all problems in the dataset {(\S\ref{subsec:problems_abstraction})}, retrieve {\bf analogous problems} {(similar to the target on an abstract level, \S\ref{subsec:similar_problems_retrieval})}, and use the top-ranked analogies and their corresponding proofs as few-shot examples 
%\dnote{we say this bit twice in the same paragraph}
{(\S\ref{subsec:llm_solver})}.
Next, we employ a {\bf symbolic verifier} that iteratively provides feedback on the validity of the generated proofs {(\S\ref{subsec:verifier})}.
%; while the proof for the target problem is unknown, we can still retrieve analogous problems -- those with similar underlying structure but different surface-level details, such as entity names or measurements.




%\noindent\textbf{(2) Verifier.} 
%We use our verifier to check the validity of the generated proof and provide feedback. This process is repeated iteratively until a correct solution is found or the LLM reaches the maximum number of retries or runs {(\S\ref{subsec:verifier})}.